% 积分图教学例：左为 4x4 原图，蓝框为第 1–2 行、第 1–2 列的矩形；右为积分图，
% 四个带符号的格子给出框内和 14-3-3+1=9。行列下标从 0 开始。宽约 68mm。
\begin{tikzpicture}[x=7mm,y=-7mm,font=\figurefont\small,
  cell/.style={draw=FigureStroke,line width=.5pt,minimum size=7mm,inner sep=0pt}]
  \fill[FigureBlueFill] (.5,.5) rectangle (2.5,2.5);
  \foreach \r/\row in {0/{1,2,0,3},1/{2,1,4,0},2/{0,3,1,2},3/{1,0,2,1}}{
    \foreach \v [count=\c from 0] in \row{\node[cell] at (\c,\r) {\v};}}
  \draw[FigureBlue,line width=1.6pt] (.5,.5) rectangle (2.5,2.5);
  \node at (1.5,-1.1) {原图 $i(x,y)$};
  \begin{scope}[xshift=40mm]
    \foreach \r/\row in {0/{1,3,3,6},1/{3,6,10,13},2/{3,9,14,19},3/{4,10,17,23}}{
      \foreach \v [count=\c from 0] in \row{\node[cell] at (\c,\r) {\v};}}
    \foreach \c/\r/\s/\f in {0/0/+/FigureBlueFill,2/0/-/FigureYellowFill,0/2/-/FigureYellowFill,2/2/+/FigureBlueFill}{
      \node[cell,fill=\f] at (\c,\r) {};
    }
    \node at (0,0) {1};\node at (2,0) {3};\node at (0,2) {3};\node at (2,2) {14};
    \foreach \c/\r/\s in {0/0/+,2/0/-,0/2/-,2/2/+}{
      \node[font=\figurefont\footnotesize,text=FigureStroke] at (\c+.32,\r-.3) {$\s$};}
    \node at (1.5,-1.1) {积分图 $ii(x,y)$};
  \end{scope}
\end{tikzpicture}
